\begin{table}[H]
\centering
\caption{Summary Statistics: Pre-Treatment Period (2014--2020)}
\label{tab:summary}
\begin{threeparttable}
\begin{tabular}{lcccc}
\toprule
 & Mean & Std.\ Dev. & Min & Max \\
\midrule
\multicolumn{5}{l}{\textit{Panel A: Privatized Municipalities (N = 
520)}} \\
  Waterborne hosp.\ rate (per 100K) & 136.3 & 187.3 & 0.5 & 1421.9 \\
  Under-5 hosp.\ rate (per 100K) & 29.5 & 40.2 & 0.0 & 325.7 \\
  Hosp.\ cost per capita (BRL) & 527.9 & 690.0 & 1.2 & 5750.6 \\
  Population & 62,878 & 362,289 & 982 & 6,747,815 \\
\midrule
\multicolumn{5}{l}{\textit{Panel B: Never-Privatized Municipalities (N = 
1,847)}} \\
  Waterborne hosp.\ rate (per 100K) & 81.8 & 129.5 & 0.0 & 1918.5 \\
  Under-5 hosp.\ rate (per 100K) & 24.4 & 36.1 & 0.0 & 454.6 \\
  Hosp.\ cost per capita (BRL) & 315.7 & 458.6 & 0.0 & 6720.7 \\
  Population & 82,680 & 442,803 & 836 & 12,325,232 \\
\bottomrule
\end{tabular}
\begin{tablenotes}[flushleft]
\small
\item \textit{Notes:} Summary statistics for the pre-treatment period (2014--2020). Privatized municipalities are those that transitioned to private water/sanitation providers following the 2020 Marco Legal do Saneamento: Alagoas Block A (2021, N = 13), CEDAE/Rio de Janeiro (2022, N = 29), Corsan/Rio Grande do Sul (2023, N = 497). Waterborne diseases defined as ICD-10 A00--A09 (intestinal infectious diseases). Data: DATASUS SIH (hospitalization records) and IBGE population estimates.
\end{tablenotes}
\end{threeparttable}
\end{table}
